Large language model inference is dominated by repeated linear projections whose computational cost and memory footprint grow rapidly with model scale. This work studies whether bit-sliced vectors (BSVs), derived from the database notion of bit-sliced indexing but used here as a numerical tensor representation, can serve as a practical compressed format for transformer inference on GPUs while preserving a native bit-sliced execution path rather than reverting to dense low-precision matrix multiplication. The work begins from a CPU-side bit-sliced vector algebra that established the first speed--accuracy--compression frontier for the representation: storage reduction was substantial, numerical fidelity could be controlled through the bit budget, but the CPU dot-product path remained too slow for transformer workloads. The later GPU implementation preserved the same representation semantics while redesigning the execution substrate around static weight conversion, dynamic query construction, and fixed-budget CUDA dot kernels. The resulting system provides a complete experimental stack that includes CPU and CUDA BSV builders, a PyTorch replacement for dense linear layers, kernel-only microbenchmarks, and end-to-end next-token evaluation on OPT-family models using LAMBADA. Under the principal GPU operating point, which uses a 7-bit query budget, a 6-bit key budget, chunk-local rescaling, and a tensor-core-oriented tile geometry, the \texttt{facebook/opt-6.7b} model decreases from \texttt{12700.031 MB} in dense form to \texttt{5022.281 MB} under full replacement while achieving \texttt{top1=0.741}, \texttt{top5=0.896}, and \texttt{66.452 ms} per-sample dot latency compared with the dense baseline at \texttt{top1=0.798}, \texttt{top5=0.929}, and \texttt{15.493 ms}. At smaller scale, however, the same representation already yields speedup on the dot path: under full replacement, \texttt{opt-125m} reduces per-sample dot latency from \texttt{3.044 ms} to \texttt{1.913 ms}. The central result is therefore not that bit-sliced execution cannot accelerate inference. The central result is that BSV compression already produces measurable speedup in some regimes, but scaling that advantage to larger models requires tighter co-design among bit budget, storage layout, runtime staging, and kernel structure.
